\section{שיעור 03 — הצגות בלתי פריקות וגנרטורים של חבורות רציפות (26.10.2025)}

\subsection*{מטרות השיעור}
\begin{itemize}
    \item הבנת מושג הפריקות והצגות בלתי פריקות
    \item למידת משפט הממדים לחבורות סופיות
    \item היכרות עם למה של שור (גרסה מורחבת) והשלכותיה הפיזיקליות
    \item הבנת הקשר בין סימטריות לניוון רמות אנרגיה
    \item מבוא לחבורות רציפות וגנרטורים אינפיניטסימליים
    \item גזירת ייצוג אקספוננציאלי של חבורות רציפות
\end{itemize}

\subsection{חזרה: סידור מחדש של איברי חבורה}

\begin{theorembox}{סידור מחדש באיברי חבורה}
בחבורה סופית $G$ עם $n$ איברים $\{g_1, g_2, \ldots, g_n\}$, מכפלת כל האיברים באיבר $g_k \in G$ מניבה סידור מחדש של אותם האיברים:
\[
\{g_k g_1, g_k g_2, \ldots, g_k g_n\} = \{g_1, g_2, \ldots, g_n\}
\]
\end{theorembox}

\textbf{הוכחה:} נובע מקיום ההופכי ומייחודיות המכפלה. אם $g_k g_i = g_k g_j$, אז בכפל מימין ב-$g_k^{-1}$ מקבלים $g_i = g_j$.

\subsection{מרחבים אינווריאנטיים ופריקות}

\subsubsection{תזכורת: מרחב אינווריאנטי}

\begin{definitionbox}{מרחב אינווריאנטי}
תת-מרחב $V' \subseteq V$ נקרא \textbf{אינווריאנטי} תחת פעולת חבורה $G$ אם:
\[
T(g) V' \subseteq V' \quad \forall g \in G
\]
כלומר, פעולת החבורה לא מוציאה וקטורים מתוך $V'$.
\end{definitionbox}

\begin{theorembox}{יצירת מרחב אינווריאנטי}
לכל וקטור $v \in V$, המרחב
\[
V' = \text{span}\{T(g)v \mid g \in G\}
\]
הוא מרחב אינווריאנטי תחת פעולת $G$.
\end{theorembox}

\subsubsection{הגדרת פריקות}

\begin{definitionbox}{הצגה פריקה ובלתי פריקה}
מרחב וקטורי $V$ עם הצגה $T$ של חבורה $G$ הוא:

\textbf{פריק (\LR{reducible}):} אם קיים תת-מרחב אינווריאנטי לא-טריוויאלי $V' \subset V$ כך ש:
\[
V = V' \oplus V'^\perp
\]
ושני המרחבים $V'$ ו-$V'^\perp$ אינווריאנטיים תחת $G$.

\textbf{בלתי פריק (\LR{irreducible, irrep}):} אם לא קיים פירוק כזה — כל תת-מרחב אינווריאנטי הוא טריוויאלי ($\{0\}$ או $V$).
\end{definitionbox}

\begin{notebox}{פרשנות פיזיקלית}
הצגה בלתי פריקה מתארת "יחידה אלמנטרית" שלא ניתן לפרק אותה עוד. בפיזיקה, אלו מתאימים למצבים עם סימטריה מוגדרת היטב.
\end{notebox}

\subsection{דוגמה מפורטת: חבורת השיקוף}

\begin{examplebox}{חבורת השיקוף $G = \{I, P\}$}
חבורה עם שני איברים: $I$ (יחידה) ו-$P$ (שיקוף).

\textbf{הצגה דו-ממדית טבעית:} במרחב $\mathbb{R}^2$ עם בסיס $\{(1,0), (0,1)\}$:
\[
I = \begin{pmatrix} 1 & 0 \\ 0 & 1 \end{pmatrix}, \quad
P = \begin{pmatrix} 0 & 1 \\ 1 & 0 \end{pmatrix}
\]

\textbf{האם ההצגה פריקה?} כן! 

\textbf{פירוק להצגות בלתי פריקות:} נבחר בסיס חדש $\{(1,1), (1,-1)\}$ (עד כדי נירמול):
\[
I = \begin{pmatrix} 1 & 0 \\ 0 & 1 \end{pmatrix}, \quad
P = \begin{pmatrix} 1 & 0 \\ 0 & -1 \end{pmatrix}
\]

שתי הצגות חד-ממדיות בלתי פריקות:
\begin{itemize}
    \item $T_1$: $I \to 1, \quad P \to 1$ (הצגה טריוויאלית)
    \item $T_2$: $I \to 1, \quad P \to -1$ (הצגת הזוגיות)
\end{itemize}
\end{examplebox}

\textbf{אנלוגיה פיזיקלית:} פונקציות זוגיות (קוסינוס) משתייכות ל-$T_1$, פונקציות אי-זוגיות (סינוס) משתייכות ל-$T_2$.

\subsection{דוגמה: חבורת התמורות $S_3$}

\subsubsection{מבנה החבורה}

\begin{definitionbox}{חבורת התמורות $S_3$}
$S_3$ היא חבורת כל התמורות של שלושה איברים.
\begin{itemize}
    \item \textbf{סדר החבורה:} $|S_3| = 3! = 6$
    \item \textbf{איברים:}
    \begin{align*}
    &e = \begin{pmatrix} 1 & 2 & 3 \\ 1 & 2 & 3 \end{pmatrix}, \quad
    (12) = \begin{pmatrix} 1 & 2 & 3 \\ 2 & 1 & 3 \end{pmatrix}, \quad
    (13) = \begin{pmatrix} 1 & 2 & 3 \\ 3 & 2 & 1 \end{pmatrix} \\
    &(23) = \begin{pmatrix} 1 & 2 & 3 \\ 1 & 3 & 2 \end{pmatrix}, \quad
    (123) = \begin{pmatrix} 1 & 2 & 3 \\ 2 & 3 & 1 \end{pmatrix}, \quad
    (132) = \begin{pmatrix} 1 & 2 & 3 \\ 3 & 1 & 2 \end{pmatrix}
    \end{align*}
\end{itemize}
\end{definitionbox}

\subsubsection{הצגה תלת-ממדית טבעית}

בבסיס $\{e_1 = (1,0,0), e_2 = (0,1,0), e_3 = (0,0,1)\}$:

\begin{examplebox}{הצגה טבעית של $S_3$}
\[
T\begin{pmatrix} 1 & 2 & 3 \\ 2 & 1 & 3 \end{pmatrix} = 
\begin{pmatrix} 0 & 1 & 0 \\ 1 & 0 & 0 \\ 0 & 0 & 1 \end{pmatrix}, \quad
T\begin{pmatrix} 1 & 2 & 3 \\ 3 & 1 & 2 \end{pmatrix} = 
\begin{pmatrix} 0 & 0 & 1 \\ 1 & 0 & 0 \\ 0 & 1 & 0 \end{pmatrix}
\]
\end{examplebox}

\textbf{האם הצגה זו פריקה?} \textbf{כן!}

\subsubsection{פירוק להצגות בלתי פריקות — קואורדינטות יעקובי}

\begin{importantbox}{פירוק ההצגה התלת-ממדית}
נבחר בסיס חדש (קואורדינטות יעקובי):
\begin{align*}
E_1 &= \frac{1}{\sqrt{3}}(1, 1, 1) \quad \text{(מרכז מאסה)} \\
E_2 &= \frac{1}{\sqrt{2}}(1, -1, 0) \quad \text{(קואורדינטה יחסית)} \\
E_3 &= \frac{1}{\sqrt{6}}(1, 1, -2) \quad \text{(קואורדינטה יחסית)}
\end{align*}

$E_1$ אינווריאנטי: $T(g) E_1 = E_1$ לכל $g \in S_3$.

המטריצות בבסיס זה מתפרקות לבלוקים:
\[
T(g) = \begin{pmatrix} 1 & \mid & 0 \\ \hline 0 & \mid & T^{(2)}(g) \end{pmatrix}
\]
כאשר $T^{(2)}(g)$ מטריצה $2 \times 2$.
\end{importantbox}

\textbf{דוגמה מפורשת:}
\begin{align*}
T'\begin{pmatrix} 1 & 2 & 3 \\ 2 & 1 & 3 \end{pmatrix} &= 
\begin{pmatrix} 1 & 0 & 0 \\ 0 & -1 & 0 \\ 0 & 0 & 1 \end{pmatrix} \\
T'\begin{pmatrix} 1 & 2 & 3 \\ 3 & 1 & 2 \end{pmatrix} &=
\begin{pmatrix} 1 & 0 & 0 \\ 0 & -\frac{1}{2} & -\frac{\sqrt{3}}{2} \\ 0 & -\frac{\sqrt{3}}{2} & \frac{1}{2} \end{pmatrix}
\end{align*}

\subsection{משפט הממדים}

\begin{theorembox}{משפט הממדים לחבורות סופיות}
תהי $G$ חבורה סופית עם $|G| = n$, ותהיינה $T_\alpha$ ($\alpha = 1, 2, \ldots, k$) כל ההצגות הבלתי פריקות הלא-שקולות של $G$, עם ממדים $d_\alpha = \dim(T_\alpha)$. אזי:
\[
\boxed{|G| = \sum_{\alpha=1}^{k} d_\alpha^2}
\]
\end{theorembox}

\textbf{משמעות:} מספר ההצגות הבלתי פריקות וממדיהן נקבעים באופן ייחודי על-ידי סדר החבורה!

\begin{examplebox}{יישום משפט הממדים}
\textbf{חבורת השיקוף:} $|G| = 2$
\[
2 = 1^2 + 1^2 \quad \Rightarrow \quad \text{שתי הצגות חד-ממדיות}
\]

\textbf{חבורת $S_3$:} $|G| = 6$
\[
6 = 1^2 + 2^2 + 1^2 \quad \Rightarrow \quad \text{שתי הצגות חד-ממדיות ואחת דו-ממדית}
\]
\end{examplebox}

\subsubsection{שלוש ההצגות הבלתי פריקות של $S_3$}

\begin{summarybox}{טבלת ההצגות של $S_3$}
\begin{center}
\begin{tabular}{|c|c|c|c|c|c|c|}
\hline
& $e$ & $(12)$ & $(13)$ & $(23)$ & $(123)$ & $(132)$ \\ \hline
$T_1$ (טריוויאלית) & $1$ & $1$ & $1$ & $1$ & $1$ & $1$ \\ \hline
$T_2$ (דו-ממדית) & $I_{2\times 2}$ & $\begin{pmatrix} -1 & 0 \\ 0 & 1 \end{pmatrix}$ & $\cdots$ & $\cdots$ & $\begin{pmatrix} -\frac{1}{2} & -\frac{\sqrt{3}}{2} \\ -\frac{\sqrt{3}}{2} & \frac{1}{2} \end{pmatrix}$ & $\cdots$ \\ \hline
$T_3$ (זוגיות) & $1$ & $-1$ & $-1$ & $-1$ & $1$ & $1$ \\ \hline
\end{tabular}
\end{center}

\textbf{$T_3$ - הצגת הזוגיות:} כל תמורה מקבלת $+1$ או $-1$ לפי זוגיות מספר ההחלפות:
\begin{itemize}
    \item תמורה \textbf{זוגית}: מספר זוגי של החלפות $\to +1$
    \item תמורה \textbf{אי-זוגית}: מספר אי-זוגי של החלפות $\to -1$
\end{itemize}
\end{summarybox}

\textbf{אימות משפט הממדים:}
\[
\dim(T_1)^2 + \dim(T_2)^2 + \dim(T_3)^2 = 1^2 + 2^2 + 1^2 = 6 = |S_3| \quad \checkmark
\]

\subsection{למה של שור (גרסה מורחבת)}

\begin{theorembox}{למה של שור 2}
תהי $G$ חבורה סופית, $T$ הצגה בלתי פריקה של $G$ במרחב $V$, ו-$C$ אופרטור הרמיטי ב-$V$.

\textbf{הנחה:} לכל $g \in G$ מתקיים:
\[
[T(g), C] = 0 \quad \text{(}C\text{ מתחלף עם כל איברי ההצגה)}
\]

\textbf{מסקנה:} $C$ פרופורציונלי לאופרטור היחידה:
\[
\boxed{C = \lambda I \quad \text{עבור } \lambda \in \mathbb{R}}
\]
\end{theorembox}

\subsubsection{הוכחת למה של שור}

\textbf{שלב 1:} $C$ הרמיטי $\Rightarrow$ קיים וקטור עצמי:
\[
Cv = \lambda v
\]

\textbf{שלב 2:} נגדיר משפחת וקטורים:
\[
\alpha = \sum_{g \in G} \alpha_g \, T(g) v
\]
כאשר $\alpha_g$ מקדמים כלשהם.

\textbf{שלב 3:} נפעיל $C$ על $\alpha$:
\begin{align*}
C \alpha &= C \sum_g \alpha_g T(g) v = \sum_g \alpha_g C T(g) v \\
&\stackrel{[C,T(g)]=0}{=} \sum_g \alpha_g T(g) C v = \sum_g \alpha_g T(g) (\lambda v) \\
&= \lambda \sum_g \alpha_g T(g) v = \lambda \alpha
\end{align*}

\textbf{מסקנה:} כל $\alpha$ מהצורה הזו הוא וקטור עצמי של $C$ עם אותו ערך עצמי $\lambda$!

\textbf{שלב 4:} המרחב
\[
V' = \text{span}\{\alpha \mid \alpha_g \in \mathbb{C}\}
\]
הוא תת-מרחב אינווריאנטי תחת $G$, וכל הווקטורים בו הם וקטורים עצמיים של $C$ עם ערך עצמי $\lambda$.

\textbf{שלב 5:} מכיוון ש-$T$ בלתי פריקה, $V'$ חייב להיות $\{0\}$ או $V$. אבל $v \in V'$, ולכן $V' \neq \{0\}$.

\textbf{מסקנה:} $V' = V$, ולכן $C$ פועל כ-$\lambda I$ על כל $V$. $\qed$

\subsection{משמעות פיזיקלית: ניוון אנרגיה}

\begin{importantbox}{קשר לפיזיקה: ניוון רמות אנרגיה}
תהי $G$ חבורת סימטריה של המילטוניאן: $[H, T(g)] = 0$ לכל $g \in G$.

\textbf{למה של שור מלמד:} אם $T_\alpha$ היא הצגה בלתי פריקה של $G$ בממד $d_\alpha$, אזי \textbf{כל $d_\alpha$ המצבים} בהצגה זו חייבים להיות בעלי \textbf{אותה אנרגיה} $E_\alpha$.

\textbf{משמעות:}
\begin{itemize}
    \item \textbf{ניוון סימטרי:} ממד ההצגה = ניוון רמת האנרגיה
    \item \textbf{ניוון מקרי:} אם שתי רמות מהצגות שונות בעלות אותה אנרגיה בדיוק (עד דיוק גבוה מאוד), יש לחשוד בסימטריה נסתרת גדולה יותר!
\end{itemize}
\end{importantbox}

\textbf{דוגמה פיזיקלית:} מולקולת מתאן (CH$_4$) בעלת סימטריית טטראדר. הספקטרום האנרגטי מראה ניוון של $1, 3, 3, 5$ — בהתאם להצגות הבלתי פריקות של חבורת הטטראדר!

\subsection{חוקי טרנספורמציה של אופרטורים}

\subsubsection{גזירת חוק הטרנספורמציה}

\textbf{מצב מקורי:} $\ket{\psi}$, אופרטור $O$, ערך תצפית:
\[
\braket{\psi | O | \psi}
\]

\textbf{מצב לאחר טרנספורמציית סימטריה:}
\[
\ket{\psi'} = U \ket{\psi}, \quad O \to O'
\]

\textbf{דרישת שקילות:} ערך התצפית נשמר:
\[
\braket{\psi | O | \psi} = \braket{\psi' | O' | \psi'}
\]

\textbf{פיתוח:}
\begin{align*}
\braket{\psi | O | \psi} &= \braket{\psi' | O' | \psi'} \\
&= \bra{\psi} U^\dagger O' U \ket{\psi}
\end{align*}

\textbf{מכיוון שזה נכון לכל $\ket{\psi}$:}
\[
\boxed{O = U^\dagger O' U \quad \Leftrightarrow \quad O' = U O U^\dagger}
\]

\textbf{פרשנות:} אופרטור עובר טרנספורמציה דמיון (conjugation) על-ידי אופרטור הסימטריה.

\subsection{מבוא לחבורות רציפות}

\subsubsection{הבדל בין חבורות דיסקרטיות לרציפות}

עד כה עסקנו בחבורות \textbf{דיסקרטיות} (בעיקר סופיות). כעת נעבור לחבורות \textbf{רציפות} — חבורות שאיבריהן תלויים בפרמטרים רציפים.

\begin{definitionbox}{חבורה רציפה}
חבורה $G$ נקראת \textbf{חבורה רציפה} אם איבריה מאופיינים על-ידי פרמטרים רציפים:
\[
g = g(\theta_1, \theta_2, \ldots, \theta_r)
\]
כאשר $(\theta_1, \ldots, \theta_r)$ משתנים רציפים על \textbf{יריעה} (\LR{manifold}) חלקה $r$-ממדית.

\textbf{דוגמאות פיזיקליות:}
\begin{itemize}
    \item \textbf{סיבובים במישור} $O(2)$: פרמטר אחד $\theta \in [0, 2\pi)$
    \item \textbf{סיבובים במרחב} $SO(3)$: שלושה פרמטרים (זוויות אוילר)
    \item \textbf{הזזות במרחב} $\mathbb{R}^3$: שלושה פרמטרים $(x, y, z)$
\end{itemize}
\end{definitionbox}

\subsection{גנרטורים אינפיניטסימליים — מפתח להבנת חבורות רציפות}

\subsubsection{דוגמה: סיבובים במישור}

נסתכל על סיבוב בזווית $\varphi$ הפועל על פונקציה $f(\theta)$ על מעגל:
\[
[T(\varphi) f](\theta) = f(\theta - \varphi)
\]

\textbf{קירוב לינארי עבור $\varphi$ קטן:}
\begin{align*}
[T(\varphi) f](\theta) &= f(\theta - \varphi) \\
&\approx f(\theta) - \varphi \frac{df}{d\theta} \\
&= \left(1 - i\varphi \cdot i\frac{d}{d\theta}\right) f(\theta)
\end{align*}

\textbf{הגדרת הגנרטור:}
\[
X = i\frac{d}{d\theta}
\]

אז:
\[
T(\varphi) \approx 1 - i\varphi X \quad \text{(עבור } \varphi \text{ קטן)}
\]

\subsubsection{גזירת הייצוג האקספוננציאלי}

\textbf{רעיון:} פירוק סיבוב גדול $\varphi$ ל-$n$ סיבובים קטנים:

\textbf{שלב 1:} סיבוב קטן:
\[
T\left(\frac{\varphi}{n}\right) \approx 1 - i\frac{\varphi}{n} X
\]

\textbf{שלב 2:} ביצוע $n$ סיבובים:
\[
T(\varphi) = \left[T\left(\frac{\varphi}{n}\right)\right]^n \approx \left(1 - i\frac{\varphi}{n} X\right)^n
\]

\textbf{שלב 3:} גבול $n \to \infty$:
\[
T(\varphi) = \lim_{n\to\infty} \left(1 - i\frac{\varphi}{n} X\right)^n = \boxed{e^{-i\varphi X}}
\]

\begin{importantbox}{הנוסחה המרכזית — ייצוג אקספוננציאלי}
\[
\boxed{T(\varphi) = e^{-i\varphi X}}
\]

כאשר:
\begin{itemize}
    \item $X = i\frac{d}{d\theta}$ הוא \textbf{גנרטור} של חבורת הסיבובים במישור
    \item $X$ הוא אופרטור הרמיטי
    \item כל איבר בחבורה מתקבל מהגנרטור והפרמטר $\varphi$
\end{itemize}
\end{importantbox}

\subsubsection{אימות באמצעות טור טיילור}

\begin{align*}
e^{-i\varphi X} f(\theta) &= e^{-i\varphi \cdot i\frac{d}{d\theta}} f(\theta) = e^{\varphi \frac{d}{d\theta}} f(\theta) \\
&= \sum_{k=0}^\infty \frac{1}{k!} \left(\varphi \frac{d}{d\theta}\right)^k f(\theta) \\
&= \sum_{k=0}^\infty \frac{\varphi^k}{k!} \frac{d^k f}{d\theta^k}(\theta) \\
&= f(\theta + \varphi) \quad \checkmark
\end{align*}

זהו בדיוק טור טיילור של $f$ סביב $\theta$!

\subsection{סיכום עיקרי הנקודות}

\begin{summarybox}{}
\textbf{הצגות בלתי פריקות:}
\begin{itemize}
    \item הצגה בלתי פריקה = "אבן בניין" יסודית שלא ניתנת לפירוק נוסף
    \item כל הצגה מתפרקת לסכום ישר של הצגות בלתי פריקות
    \item פירוק על-ידי מעבר לבסיס מתאים (לכסון בלוקים)
\end{itemize}

\textbf{משפט הממדים:}
\[
|G| = \sum_{\alpha} d_\alpha^2
\]
\begin{itemize}
    \item סדר החבורה קובע את מספר וממדי ההצגות הבלתי פריקות
    \item כלי חזק לזיהוי כל ההצגות האפשריות
\end{itemize}

\textbf{למה של שור 2:}
\begin{itemize}
    \item אם $C$ הרמיטי ומתחלף עם כל $T(g)$ בהצגה בלתי פריקה, אז $C = \lambda I$
    \item שימוש פיזיקלי: $C = H$ (המילטוניאן) $\Rightarrow$ ניוון = ממד ההצגה
\end{itemize}

\textbf{משמעות פיזיקלית — ניוון אנרגיה:}
\begin{itemize}
    \item סימטריה של המילטוניאן $\Rightarrow$ ניוון רמות אנרגיה
    \item ממד ההצגה הבלתי פריקה = ניוון רמת האנרגיה
    \item ניוון מקרי (לא סימטרי) מצביע על סימטריה נסתרת
\end{itemize}

\textbf{חוקי טרנספורמציה:}
\[
O' = U O U^\dagger
\]
אופרטור עובר טרנספורמציה דמיון תחת סימטריה.

\textbf{חבורות רציפות:}
\begin{itemize}
    \item איברים מאופיינים על-ידי פרמטרים רציפים על יריעה
    \item גנרטורים אינפיניטסימליים $X$ "חיים" בשכונת איבר היחידה
    \item ייצוג אקספוננציאלי: $T(\varphi) = e^{-i\varphi X}$
    \item גנרטורים הם אופרטורים הרמיטיים
\end{itemize}
\end{summarybox}

\subsection*{נקודות למחשבה ובדיקת הבנה}

\begin{enumerate}
    \item הסבירו את ההבדל בין הצגה פריקה להצגה בלתי פריקה. מדוע ההצגות הבלתי פריקות חשובות יותר?
    
    \item עבור חבורה בסדר $|G| = 12$, מהם הפירוקים האפשריים לפי משפט הממדים? (רמז: חפשו פתרונות ל-$12 = \sum d_\alpha^2$)
    
    \item הסבירו מדוע ניוון של רמת אנרגיה במערכת פיזיקלית קשור לממד ההצגה הבלתי פריקה של חבורת הסימטריה.
    
    \item מדוע הגנרטור $X = i\frac{d}{d\theta}$ של חבורת הסיבובים חייב להיות הרמיטי?
    
    \item הוכיחו ש-$e^{\alpha \frac{d}{dx}} f(x) = f(x + \alpha)$ באמצעות טור טיילור.
    
    \item תנו דוגמה פיזיקלית למערכת שבה ניוון האנרגיה נקבע על-ידי סימטריה. איזו חבורה רלוונטית?
\end{enumerate}

\newpage
